% Springer LNCS proceedings paper (MICCAI, ECML PKDD and other Springer conferences).
% Uses Springer's llncs class and splncs04 bibliography style, as distributed on CTAN.
% Check the conference's call for papers for its page limit and anonymisation rules.
\documentclass[runningheads]{llncs}

\usepackage[T1]{fontenc}
\usepackage{graphicx}
\usepackage{amsmath,amssymb}
\usepackage{booktabs}
% Coloured links help on screen; Springer removes them in print.
\usepackage[colorlinks=true,allcolors=blue]{hyperref}

\begin{document}

\title{Contribution Title\thanks{Supported by organization x.}}
\titlerunning{Abbreviated Title}

% For a double-blind submission, replace the authors with "Anonymous Authors".
\author{First Author\inst{1}\orcidID{0000-1111-2222-3333} \and
Second Author\inst{2,3}\orcidID{1111-2222-3333-4444} \and
Third Author\inst{3}\orcidID{2222-3333-4444-5555}}
\authorrunning{F. Author et al.}

\institute{Princeton University, Princeton NJ 08544, USA \and
Springer Heidelberg, Tiergartenstr. 17, 69121 Heidelberg, Germany
\email{lncs@springer.com}\\
\url{https://www.springer.com/gp/computer-science/lncs} \and
ABC Institute, Rupert-Karls-University Heidelberg, Heidelberg, Germany\\
\email{\{abc,lncs\}@uni-heidelberg.de}}

\maketitle

\begin{abstract}
The abstract should briefly summarize the contents of the paper in 150--250 words.

\keywords{First keyword \and Second keyword \and Another keyword.}
\end{abstract}

\section{Introduction}
Springer's LNCS format sets papers in a single column with running heads. Headings of the
first and second level are numbered; the third level is run in, as below.

\subsubsection{A run-in heading.}
Run-in headings are followed by the text on the same line.

\section{Method}
Equations are numbered on the right:
\begin{equation}
  \mathcal{L}(\theta) = \frac{1}{N} \sum_{i=1}^{N} \ell\bigl(f_\theta(x_i), y_i\bigr).
\end{equation}

\begin{theorem}
The class provides theorem-like environments such as theorem, lemma, proposition, definition,
corollary, example and remark.
\end{theorem}

\begin{proof}
Proofs, examples and remarks are set in plain type.
\end{proof}

\section{Experiments}
Table captions go above the table (Table~\ref{tab:results}); figure captions go below the
figure.

\begin{table}
  \caption{Table captions should be placed above the tables.}
  \label{tab:results}
  \centering
  \begin{tabular}{lcc}
    \toprule
    Method & Accuracy (\%) & Time (s) \\
    \midrule
    Baseline & 81.2 & 12.4 \\
    Ours & \textbf{85.7} & 10.1 \\
    \bottomrule
  \end{tabular}
\end{table}

\section{Conclusion}
For citations of references, we prefer the use of square brackets and consecutive numbers,
e.g.\ \cite{ref_article1,ref_lncs1,ref_book1}.

\begin{credits}
\subsubsection{\ackname} A bold run-in heading in small font size at the end of the paper is
used for general acknowledgments.

\subsubsection{\discintname}
The authors have no competing interests to declare that are relevant to the content of this
article.
\end{credits}

% Springer asks for splncs04 with BibTeX; the references are inline here so the sample builds
% without a .bib file.
\begin{thebibliography}{8}
\bibitem{ref_article1}
Author, F.: Article title. Journal \textbf{2}(5), 99--110 (2016)

\bibitem{ref_lncs1}
Author, F., Author, S.: Title of a proceedings paper. In: Editor, F., Editor, S. (eds.)
CONFERENCE 2016, LNCS, vol. 9999, pp. 1--13. Springer, Heidelberg (2016).
\doi{10.10007/1234567890}

\bibitem{ref_book1}
Author, F., Author, S., Author, T.: Book title. 2nd edn. Publisher, Location (1999)
\end{thebibliography}

\end{document}
